\section{Introduction}

Parameter-efficient fine-tuning methods like LoRA \cite{hu2021lora} have revolutionized adaptation of large language models, enabling task-specific customization with minimal compute. Yet their success relies on an implicit assumption: that the base model architecture already supports the target task. When this assumption breaks, even the most sophisticated fine-tuning fails catastrophically. A state-space model trained for generation can achieve 81\% accuracy on sentiment classification but performs 12 percentage points \textit{below} random baseline on paraphrase detection---a failure no amount of parameter tuning can overcome.

This phenomenon reveals a fundamental gap in our understanding of parameter-efficient fine-tuning across diverse architectures. While LoRA and related methods \cite{li2021prefix,lester2021power} have been extensively validated on transformer architectures, their applicability to emerging sub-quadratic models like state-space models \cite{gu2023mamba} remains unexplored. Sub-quadratic architectures promise efficiency gains through linear-time processing, replacing the quadratic attention mechanism with recurrent state-space dynamics. However, these architectural differences introduce constraints: causal state-space models process text autoregressively without bidirectional context, fundamentally limiting which tasks they can perform.

The core challenge is \textit{architectural compatibility}: some tasks require capabilities inherent to the model architecture, independent of learned parameters. Paraphrase detection requires symmetric comparison of two text sequences---a capability present in transformer encoders through bidirectional attention but absent in causal decoders. If the base architecture cannot perform a task at zero-shot evaluation, parameter-efficient fine-tuning cannot add the missing capability. LoRA adapts existing representations but cannot fundamentally alter the information flow dictated by model architecture.

We demonstrate this principle through systematic evaluation of LoRA transfer from transformers to state-space models. Our key insight is that \textbf{zero-shot evaluation serves as a critical architectural compatibility gate}: if a pretrained model fails a task before fine-tuning, the architectural mismatch signals that PEFT will fail regardless of hyperparameter tuning. This simple validation step prevents wasted compute on structurally impossible experiments.

We evaluate Mamba-130M \cite{gu2023mamba}, a representative causal state-space model, on three GLUE classification tasks \cite{wang2018glue}: natural language inference (MNLI), paraphrase detection (QQP), and sentiment analysis (SST-2). These tasks span different reasoning requirements: MNLI and QQP require bidirectional comparison of premise-hypothesis or question-question pairs, while SST-2 aligns with causal language modeling's left-to-right structure. Zero-shot evaluation reveals a striking binary pattern: Mamba achieves 81\% accuracy on SST-2 (+31pp above random, $p<0.001$), demonstrating genuine language understanding. Yet it achieves only 38\% on QQP ($-12$pp below random, $p=0.003$), performing worse than guessing. This below-random performance on paraphrase detection exposes architectural incompatibility that no fine-tuning strategy can overcome.

Building on this insight, we make the following contributions:

\begin{enumerate}
\item \textbf{Systematic study of PEFT applicability to state-space models}: We evaluate whether parameter-efficient fine-tuning methods developed for transformers apply to causal state-space models, revealing that PEFT assumes architectural alignment with target tasks.

\item \textbf{Empirical demonstration of architectural failure modes}: Mamba-130M fails bidirectional reasoning tasks at zero-shot (QQP: 38\% vs 50\% baseline), confirming that causal architectures cannot perform symmetric comparisons required for paraphrase detection.

\item \textbf{Validation of zero-shot evaluation as architectural gate}: We show that zero-shot performance reliably predicts PEFT viability, preventing wasted compute on architecturally incompatible experiments.

\item \textbf{Task-architecture compatibility framework}: We establish that architectural capabilities are prerequisites for PEFT, not outcomes. Fine-tuning specializes existing capabilities but cannot add fundamentally new ones.
\end{enumerate}

\textbf{Scope and Limitations:} We evaluate zero-shot architectural compatibility without empirically testing LoRA fine-tuning after failure. Our hypothesis predicts fine-tuning would amplify architectural bias rather than overcome it, but empirical validation remains future work. Results are based on a single SSM instance (Mamba-130M); broader claims about causal SSMs require testing additional variants.

Our findings challenge the implicit assumption in PEFT literature that methods developed for transformers transfer seamlessly to other architectures. As sub-quadratic models proliferate to address transformer scalability limitations, practitioners must validate architectural compatibility before investing in fine-tuning infrastructure. Zero-shot evaluation provides a principled, low-cost method to identify doomed experiments early in the research pipeline.
